\ExamPart{Câu tự luận}{3}{ }

\NQuestion Giải các yêu cầu sau:
\begin{SubQuestions}
  \item Giải bất phương trình $\left(x^2-7x+10\right)\left(x^2-4\right)\le0$.
  \item Tìm tất cả các giá trị của $m$ để bất phương trình $x^2-2mx+m+2>0$ đúng với mọi $x\in\mathbb{R}$.
\end{SubQuestions}
\AnswerLine{$x\in[-2;5],\ m\in(-1;2)$}
\SolutionBlock{\begin{SubQuestions}
\item Ta có
\[
x^2-7x+10=(x-2)(x-5),\qquad x^2-4=(x-2)(x+2).
\]
Suy ra
\[
\left(x^2-7x+10\right)\left(x^2-4\right)=(x-2)^2(x-5)(x+2).
\]
Vì $(x-2)^2\ge0$ với mọi $x$, nên dấu của biểu thức phụ thuộc vào $(x-5)(x+2)$. Do đó
\[
(x-2)^2(x-5)(x+2)\le0 \iff (x-5)(x+2)\le0.
\]
Suy ra
\[
-2\le x\le 5.
\]
Vậy tập nghiệm là $[-2;5]$.
\item Để bất phương trình
\[
x^2-2mx+m+2>0
\]
đúng với mọi $x\in\mathbb{R}$, cần và đủ:
\[
a=1>0\quad \text{và}\quad \Delta<0.
\]
Ta có
\[
\Delta=(-2m)^2-4(m+2)=4m^2-4m-8=4(m^2-m-2)=4(m-2)(m+1).
\]
Điều kiện $\Delta<0$ tương đương
\[
(m-2)(m+1)<0\iff -1<m<2.
\]
Vậy $m\in(-1;2)$.
\end{SubQuestions}}

\NQuestion Giải các phương trình sau:
\begin{SubQuestions}
  \item $x^2-5x+6=0$.
  \item $x^4-10x^2+9=0$.
\end{SubQuestions}
\AnswerLine{$x=2,3$; $x=\pm1,\pm3$}
\SolutionBlock{\begin{SubQuestions}
\item Ta có
\[
x^2-5x+6=(x-2)(x-3)=0.
\]
Vậy nghiệm là $x=2$ hoặc $x=3$.
\item Đặt $t=x^2\ (t\ge0)$, ta được
\[
t^2-10t+9=0\iff (t-1)(t-9)=0.
\]
Suy ra $t=1$ hoặc $t=9$, nên
\[
x=\pm1,\ \pm3.
\]
Vậy phương trình có bốn nghiệm $-3,-1,1,3$.
\end{SubQuestions}}

\NQuestion Trong mặt phẳng tọa độ $Oxy$, cho ba điểm $A(1;2)$, $B(5;4)$, $C(3;-2)$.
\begin{SubQuestions}
  \item Tính tọa độ các vectơ $\overrightarrow{AB}$ và $\overrightarrow{AC}$.
  \item Chứng minh rằng tam giác $ABC$ vuông tại $A$.
  \item Viết phương trình đường thẳng $BC$.
  \item Tính khoảng cách từ điểm $A$ đến đường thẳng $BC$.
\end{SubQuestions}
\AnswerLine{$\overrightarrow{AB}=(4;2),\ \overrightarrow{AC}=(2;-4),\ BC:3x-y-11=0,\ d(A,BC)=\dfrac{10}{\sqrt{10}}=\sqrt{10}$}
\SolutionBlock{\begin{SubQuestions}
\item Ta có
\[
\overrightarrow{AB}=(5-1;4-2)=(4;2),\qquad \overrightarrow{AC}=(3-1;-2-2)=(2;-4).
\]
\item Tính tích vô hướng:
\[
\overrightarrow{AB}\cdot\overrightarrow{AC}=4\cdot2+2\cdot(-4)=8-8=0.
\]
Suy ra $\overrightarrow{AB}\perp\overrightarrow{AC}$, nên tam giác $ABC$ vuông tại $A$.
\item Đường thẳng $BC$ đi qua $B(5;4)$ và $C(3;-2)$ nên có hệ số góc
\[
m=\dfrac{-2-4}{3-5}=3.
\]
Phương trình qua $B$ là
\[
y-4=3(x-5)\iff 3x-y-11=0.
\]
\item Khoảng cách từ $A(1;2)$ đến đường thẳng $BC$ là
\[
d(A,BC)=\dfrac{|3\cdot1-2-11|}{\sqrt{3^2+(-1)^2}}
=\dfrac{10}{\sqrt{10}}=\sqrt{10}.
\]
\end{SubQuestions}}

\NQuestion Trong mặt phẳng tọa độ $Oxy$, cho đoạn thẳng $AB$ với $A(-1;3)$, $B(5;1)$.
\begin{SubQuestions}
  \item Tìm tọa độ trung điểm $M$ của đoạn thẳng $AB$.
  \item Viết phương trình đường trung trực của đoạn thẳng $AB$.
  \item Gọi $d$ là đường thẳng đi qua $A$ và vuông góc với $AB$. Chứng minh rằng $d$ song song với đường trung trực của $AB$.
\end{SubQuestions}
\AnswerLine{$M(2;2)$,\ $\Delta:3x-y-4=0$}
\SolutionBlock{\begin{SubQuestions}
\item Trung điểm của $AB$ là
\[
M\left(\dfrac{-1+5}{2};\dfrac{3+1}{2}\right)=(2;2).
\]
\item Ta có
\[
\overrightarrow{AB}=(5-(-1);1-3)=(6;-2)=(3;-1)\cdot2.
\]
Vì đường trung trực của $AB$ vuông góc với $AB$ nên nhận một vectơ pháp tuyến cùng phương với $\overrightarrow{AB}$, chẳng hạn $(3;-1)$. Đi qua $M(2;2)$, nên phương trình là
\[
3(x-2)-(y-2)=0\iff 3x-y-4=0.
\]
\item Đường thẳng $d$ đi qua $A$ và vuông góc với $AB$, nên $d$ cũng có vectơ pháp tuyến cùng phương với $\overrightarrow{AB}=(3;-1)$. Do đó $d$ và đường trung trực của $AB$ đều vuông góc với $AB$, suy ra chúng song song với nhau.
\end{SubQuestions}}
